\documentclass[12pt]{article}

% --- INICIO DE TU CONFIGURACIÓN NECESARIA ---
\usepackage[spanish, es-tabla, es-notilde]{babel}
\usepackage[utf8]{inputenc}
\usepackage[a4paper,top=3cm,bottom=2.5cm,left=3cm,right=3cm]{geometry}
\usepackage{mathtools, amsfonts, amssymb, mathrsfs}

% --- PAQUETES NUEVOS PARA EL ALGORITMO ---
\usepackage{algorithm}
\usepackage{algorithmic}

% Para que el título del algoritmo aparezca en español ("Algoritmo" en lugar de "Algorithm")
\floatname{algorithm}{Algoritmo}
\input{../../_assets/preambulo.tex}

\begin{document}

    % 1. Foto de fondo
    % 2. Título
    % 3. Encabezado Izquierdo
    % 4. Color de fondo
    % 5. Coord x del titulo
    % 6. Coord y del titulo
    % 7. Fecha

    
    \input{../../_assets/portada}
    \portadaExamen{ffccA4.jpg}{Modelos \\ Avanzados de \\ Computación \\Examen II \\ de Prácticas}{Modelos Avanzados de Computación. Examen II de Prácticas}{MidnightBlue}{-8}{28}{2026}{José Juan Urrutia Milán}

    \newpage

    \begin{description}
        \item[Asignatura] Modelos Avanzados de Computación.
        \item[Curso Académico] 2025/26.
        \item[Grado] Grado en Ingeniería Informática.
        \item[Grupo] Grupo de Prácticas A1.
        \item[Profesor] Serafín Moral Callejón.
        \item[Descripción] Tercer Examen de Prácticas.
        \item[Fecha] 27 de Mayo.
        \item[Duración] 1 hora.
    \end{description}
    \newpage


    % ------------------------------------

    \begin{ejercicio}
        Entre los dos siguientes problemas uno es espacio-logarítmico y el otro es tiempo-polinómico. Clasificar los problemas y dar un algoritmo para demostrar sus clases:
        \begin{enumerate}
            \item Dado un conjunto de números, determinar la mediana del conjunto.

                En caso de que la cantidad de números no sea impar, rechazar.
            \item Dado un conjunto de números, determinar si existe una asignación en parejas de modo que la suma de todas las parejas sea la misma.
        \end{enumerate}
    \end{ejercicio}

    \newpage
    \setcounter{ejercicio}{0}
    \noindent
    \textbf{Solución.}

    \begin{ejercicio}
        Entre los dos siguientes problemas uno es espacio-logarítmico y el otro es tiempo-polinómico. Clasificar los problemas y dar un algoritmo para demostrar sus clases:
        \begin{enumerate}
            \item Dado un conjunto de números, determinar la mediana del conjunto.

                En caso de que la cantidad de números no sea impar, rechazar.
            \item Dado un conjunto de números, determinar si existe una asignación en parejas de modo que la suma de todas las parejas sea la misma.
        \end{enumerate}

        \noindent
        Podemos intuir que el primer problema es más fácil que el segundo, por lo que trataremos de crear un algoritmo espacio-logarítmico para el primero y un algoritmo tiempo-polinómico para el segundo. A pesar de ser el segundo problema de una clase más difícil, resulta más fácil la tarea de crear dicho algoritmo.
        \begin{description}
            \item [Problema 2.] Debemos crear un algoritmo tiempo-polinómico para esta tarea. El algoritmo es el siguiente:

            \begin{algorithm}
                \caption{}\label{alg:1}
            \begin{algorithmic}
                \STATE \textbf{Entrada:} Una secuencia de números $n_1, n_2, \ldots, n_m$ separados por $\&$. \\
                \noindent
                \STATE Contar la cantidad $m$ de números en una segunda cinta.
                \STATE \textbf{Si} $m$ es impar \textbf{entonces}
                \STATE \quad Terminar y decir ``NO''.
                \STATE \textbf{Fin Si} 
                \STATE Ordenar los números de menor a mayor.
                \STATE Almacenar en una tercera cinta la suma del primero y del último, $k$.
                \STATE Poner ``2'' en una cuarta cinta, nos referimos a este contador por $i$.
                \STATE \textbf{Mientras que} $i\leq \nicefrac{m}{2}$ \textbf{entonces}
                \STATE \quad $t = $ Sumar el elemento $i$ más el $m-i+1-$ésimo.
                \STATE \quad \textbf{Si} $t \neq k$ \textbf{entonces}
                \STATE \quad \quad Terminar y decir ``NO''.
                \STATE \quad \textbf{Fin Si}
                \STATE \quad Incrementar $i$ en una unidad.
                \STATE \textbf{Fin Mientras}
                \STATE Terminar y decir ``SÍ''.
            \end{algorithmic}
            \end{algorithm}

            Vemos a continuación:

            \begin{itemize}
                \item El Algoritmo~\ref{alg:1} es tiempo-polinómico. Supuesto que el número de caracteres de entrada $0,1,\&$ es $n$:
                    \begin{enumerate}
                        \item Contar la cantidad $m$ de números es $O(n)$.
                        \item Ordenar los números de menor a mayor puede hacerse en\footnote{Sabemos que puede hacerse en $O(n\log n)$, pero esta suposición facilita el razonamiento.} $O(n^2)$.
                        \item El procedimiento del ``mientras'' recorre $m-2$ casillas, por lo que es $O(n)$, y su cuerpo es constante.
                    \end{enumerate}
                    En definitiva, el algoritmos es $O(n^2)$, por lo que es tiempo-polinómico.
                \item El Algoritmo~\ref{alg:1} es correcto (es decir, resuelve el problema), puesto que si existe una tal asignación, el número máximo $M$ debe pertenecer a alguna pareja, por lo que debe ser emparejado con otro número $n_k$. Si este número no es el más pequeño de todos ($m < n_k$), al emparejar el más pequeño con otro número $n_t$ (tenemos $n_t\leq M$) tendremos que:
                    \begin{equation*}
                        n_t + m < n_k + M
                    \end{equation*}
                    Por lo que este emparejamiento no sería válido. Así, si un emparejamiento de este tipo existe, el número máximo debe ser emparejado con el mínimo. Si ahora consideramos el mismo problema pero sin considerar los números máximo y mínimo, por el mismo razonamiento llegamos a que el segundo más grande debe ser emparejado con el segundo más pequeño.

                    De esta forma, se puede demostrar por inducción que, en caso de existir una tal asignación, debe ser la que hemos hecho.
            \end{itemize}

            \item [Problema 1.] Buscamos un algoritmo espacio-logarítmico:

                \begin{algorithm}
                    \caption{}\label{alg:2}
                \begin{algorithmic}
                    \STATE \textbf{Entrada:} Una secuencia de números en binario $n_1,n_2,\ldots,n_m$ separados por \&.\\

                    \noindent
                    \STATE Contar la cantidad de números $m$ en una segunda cinta con un contador binario.
                    \STATE \textbf{Si} $m$ es par \textbf{entonces}.
                    \STATE \quad Terminar y decir ``NO''.
                    \STATE \textbf{Fin Si}
                    \STATE Copiar el primer número $n_1$ en una tercera cinta, separado por \& de su índice.
                    \STATE Inicializar la cuarta cinta con un contador binario $c$ inicializado a $0$.
                    \STATE \textbf{Repetir} $m$ veces:
                    \STATE \quad Usar $c$ para contar la cantidad de números menores o iguales que el de la tercera cinta. 
                    \STATE \quad Usar una quinta cinta para guardar el número $n_k$ que consigue un menor valor de $c$ siendo $c>\frac{m-1}{2}$.
                    \STATE \quad \textbf{Si} $c = \frac{m-1}{2}$ \textbf{entonces}
                    \STATE \quad \quad Terminar y devolver el número de la tercera cinta.
                    \STATE \quad \textbf{Fin Si}
                    \STATE \quad Si en la tercera cinta tenemos ``$n_i\& i$'', cambiarlo por $n_{i+1}\&i+1$.
                    \STATE \textbf{Fin Repetir}
                    \STATE Terminar y devolver el número de la quinta cinta.
                \end{algorithmic}
                \end{algorithm}

                \noindent
                Vemos:
                \begin{itemize}
                    \item El algoritmo es espacio-logarítmico. Para verlo, veamos qué espacio usamos, suponiendo que la secuencia de entrada contiene $n$ caracteres:
                        \begin{itemize}
                            \item La secuencia de números de entrada nunca se modifica, y la secuencia de salida la escribimos de izquierda a derecha sin volver hacia atrás, por lo que estas secuencias no las contamos.
                            \item En la segunda cinta se cuenta el número $m$ de números en binario. Contar en binario hasta $n$ ocupa un espacio $O(\log n)$, por lo que como $m<n$ tendremos que este contador es $O(\log n)$.
                            \item En la tercera guardamos números de la secuencia de entrada separados por su índice en binario, por lo que también es $O(\log n)$.
                            \item En la cuarta cinta guardamos un contador binario que siempre será menor o igual que $m$, por lo que es $O(\log n)$.
                        \end{itemize}
                        En definitiva, el algoritmo es $O(\log n)$ en espacio.
                    \item El Algoritmo~\ref{alg:2} es correcto, puesto que si $m$ es impar siempre tendremos que $\exists k \in \{1,\ldots,m\}$ con $n_k$ siendo la mediana, por lo que:
                        \begin{itemize}
                            \item Bien $c = \frac{m-1}{2}$ para $n_k$, de donde $n_k$ es la mediana.
                            \item Bien $c$ para $n_k$ es estrictamente mayor que $\frac{m-1}{2}$ pero es el mínimo valor de $c$ para cualquier otro valor $n_t\neq n_k$. Esta situación ocurre cuando el valor de la mediana se repite, como por ejemplo en:
                                \begin{equation*}
                                    3\ 1\ 4\ 4\ 4\ 5\ 7
                                \end{equation*}
                                Tenemos que $m = 7\Longrightarrow \frac{m-1}{2}= 3$. Vemos que $c=2$ para $n_1$, que $c=1$ para $n_2$, $c=5$ para $n_3=n_4=n_5$, $c=6$ para $n_6$ y $c=7$ para $n_7$.

                                En este caso la mediana se alcanza en $4$, cuyo valor $c$ es mayor que $\frac{m-1}{2}$ pero es el mínimo con esta propiedad.
                        \end{itemize}
                \end{itemize}

        \end{description}
    \end{ejercicio}

\end{document}
